% National Institute of Technology, Srinagar Thesis Template for LaTeX 

% (unofficial) Thesis Template for LaTeX

% Edited by Nayeem Iqbal


% DON'T USE THIS TEMPLATE IF YOU DON'T KNOW WHAT YOU'RE DOING!
% Remember, it comes WITH NO WARRANTY!


% Here is how to use this template:
%
% DON'T FORGET TO ADD YOUR OWN NAME AND TITLE and other changes in the certificates, similarity index, no of pages

% The template is based on the standard "book" document class which provides all necessary sectioning structures and allows multi-part theses.

% Tip 1: Generate your figures and photos in the size you want them to appear
% in your thesis, rather than scaling them with \includegraphics options.
% Tip 2: Any drawings you do should be in scalable vector graphic formats:
% SVG, PNG, WMF, EPS and then converted to PNG or PDF, so they are scalable in
% the final PDF as well.
% Tip 3: Photographs should be cropped and compressed so as not to be too large.

% To create a PDF output that is optimized for double-sided printing: 
%
% 1) comment-out the \documentclass statement in the preamble below, and
% un-comment the second \documentclass line.
%
% 2) change the value assigned below to the boolean variable
% "PrintVersion" from "false" to "true".

% --------------------- Start of Document Preamble -----------------------

% Specify the document class, default style attributes, and page dimensions
% For hyperlinked PDF, suitable for viewing on a computer, use this:
\documentclass[a4paper,12pt,titlepage,oneside,final]{book}
 
%our college require one side printing, for thesis submission.
 
 
% For PDF, suitable for double-sided printing, change the PrintVersion variable below
% to "true" and use this \documentclass line instead of the one above:
% \documentclass[a4paper,12pt,titlepage,openright,twoside,final]{book}


% This package allows if-then-else control structures.
\usepackage{ifthen}
\newboolean{PrintVersion}
\setboolean{PrintVersion}{false} 
% \setboolean{PrintVersion}{true} 
% CHANGE THIS VALUE TO "true" as necessary, to improve printed results 
% for hard copies by overriding some options of the hyperref package.


% Load your needed packages and other commands of yours.
\input{thesis-preamble}

% This is where thesis margins and spaces are set.
\input{private/thesis-margins-and-spaces}

%======================================================================
%   L O G I C A L    D O C U M E N T -- the content of your thesis
%======================================================================
\begin{document}

% For a large document, it is a good idea to divide your thesis
% into several files, each one containing one chapter.
% To illustrate this idea, the "front pages" (i.e., title page,
% declaration, certificate, similarity index, abstract, acknowledgements,
% dedication, table of contents, list of tables, list of figures, list of
% abbrevations, list of symbols etc.
%----------------------------------------------------------------------
% FRONT MATERIAL
%----------------------------------------------------------------------
%
% C O V E R  P A G E
% ------------------
\newcommand{\thesisauthor}{Student Name}
\newcommand{\thesistitlecoverpage}{NIT Srinagar thesis format in {LaTeX}}
\newcommand{\nameofdegree}{Master of Technology} % Bachelor of technology, MSc, Phd, MS, MBA etc use full form here..
\newcommand{\nameofprogram}{Micro Electronics} 
\newcommand{\nameofguide}{Guide Name} %without prefix
\newcommand{\enrol}{202XMECEMEXXX} 
\newcommand{\academicunit}{Department of Electronics and Communication Engineering}
\newcommand{\graduationyear}{July 2022} %edit to your own graduation year
%
\input{private/frontpage}
%
%
% R E S T  O F  F R O N T  P A G E S
% ----------------------------------
\input{private/certificate}
\input{private/similarity_certificate} %minimal edit with your similarity index
\input{private/declarationpage}

% Edit the following 3 files with your abstract, acknowledgements, 
% and dedication.
\input{Chapters/Abstract}
\input{Chapters/Acknowledgements}
\input{Chapters/Dedication} %You can comment-out dedication, if you want. This is not necessary.

% No need to edit this file.
\input{private/toc-lot-lof}

%List the abbrevations. You can comment-out List_of_Abbreviations and symbols, if you dont want to have abbrevations and symbols
\input{Chapters/List_of_Abbreviations}
\input{Chapters/List_of_Symbols} 



% Change page numbering back to Arabic numerals
\pagenumbering{arabic}

%----------------------------------------------------------------------
% MAIN BODY
%---------------------------------------------------------------------- 
% Chapters 
% Include your "sub" source files here (must have extension .tex)
\include{Chapters/Chapter_01}
\include{Chapters/Chapter_02}
\include{Chapters/Chapter_03}

%Add Chapters as much as you want!
%3 chapters are suggested
%1st introduction and review
%2nd physics and problem statement
%3rd result, conclusion and future outlook
%remain in consult with your advisor about the no of chapters and the structure...



% B I B L I O G R A P H Y
% -----------------------
%
% The following statement selects the style to use for references.  It controls the sort order of the entries in the bibliography and also the formatting for the in-text labels.
\bibliographystyle{plainurl}
% This specifies the location of the file containing the bibliographic information.  
% It assumes you're using BibTeX (if not, why not?).



\ifthenelse{\boolean{PrintVersion}}{
\cleardoublepage % This is needed if the book class is used, to place the anchor in the correct page,
                 % because the bibliography will start on its own page.
}{
\clearpage       % Use \clearpage instead if the document class uses the "oneside" argument
}


\bibliography{bibliography/references}
% Add the References to the Table of Contents
\phantomsection 
\addcontentsline{toc}{chapter}{\textbf{References}}
% Tip 5: You can create multiple .bib files to organize your references. 
% Just list them all in the \bibliogaphy command, separated by commas (no spaces).
\phantomsection 

\begin{appendix}
\include{Chapters/Appendix_A} %"Software & Code for obtaining the plots and results"

%include appendix as required, you can also include proof of concept.

\end{appendix}


%----------------------------------------------------------------------
% END MATERIAL
%----------------------------------------------------------------------

%----------------------------------------------------------------------
\end{document}
%======================================================================


